\documentclass[11pt,a4paper]{article}

\usepackage[T1]{fontenc}
\usepackage[utf8]{inputenc}
\usepackage[brazil]{babel}
\usepackage{lmodern}
\usepackage{microtype}
\usepackage[margin=2.5cm]{geometry}
\usepackage{booktabs}
\usepackage{array}
\usepackage{enumitem}
\usepackage{xcolor}
\usepackage{listings}
\usepackage{hyperref}
\usepackage{url}

\hypersetup{
  colorlinks=true,
  linkcolor=blue!55!black,
  urlcolor=blue!55!black,
  citecolor=blue!55!black,
  pdftitle={Estado causal persistente para enforcement no sistema operacional}
}

\title{Estado Causal Persistente para Enforcement no Sistema Operacional\\
\large Engenharia de um Runtime de Política Linux com Estado, Auditoria e Falha Fechada}

\author{
Felipe Maya Muniz\\
\textit{Unix-AGB Project / AletheionAGI}
}

\date{17 de agosto de 2026}

\begin{document}
\maketitle

\begin{abstract}
Mecanismos tradicionais de enforcement no sistema operacional decidem a partir
da requisição atual: executável, credencial, caminho, endereço de rede ou rótulo
estático. Essa visão local é eficiente e confiável, mas não expressa diretamente
um risco que emerge de uma sequência de ações individualmente comuns. O Unix-AGB
(Aletheion Guard Bridge) investiga uma arquitetura complementar na qual eventos
normalizados do kernel atualizam estado causal persistente e isolado por
namespace; uma camada determinística de política compila apenas decisões
restritivas em registros de enforcement de curta duração. Eventos exatos
permanecem como verdade histórica, estado aprendido continua substituível e o
Linux preserva a autoridade final de enforcement. A restrição central do
projeto é a \emph{autoridade monotônica}: inferência aprendida ou com estado só
pode reduzir privilégios existentes; não pode fabricar permissão. Este artigo
explica a engenharia desse pipeline: contratos, identidade, ordenação, persistência,
revisões de política, caches autenticados, handoff por seccomp user notification,
recuperação, empacotamento e validação experimental. Também delimita o resultado:
o sistema atual é um protótipo controlado com enforcement de egress empacotado,
opt-in e limitado a processos lançados pelo mecanismo; não é um serviço de
segurança de produção e ainda não demonstra eficácia contra ataques naturais
desconhecidos.
\end{abstract}

\noindent\textbf{Palavras-chave:}
estado causal persistente,
segurança de sistemas operacionais,
aplicação de políticas com estado,
autoridade monotônica,
compilação somente de negações,
segurança Linux,
seccomp,
segurança em tempo de execução.

\section{O problema de engenharia}

Considere dois processos realizando a mesma operação terminal: abrir uma
credencial ou conectar a um endereço externo. Um pode ter seguido um fluxo
normal de configuração; o outro pode ter executado uma ferramenta incomum,
estabelecido conectividade externa e só então acessado o recurso sensível. Uma
regra sem estado vê requisições terminais idênticas. Um sistema com estado
consegue distinguir os históricos, mas introduzir estado na política do sistema
operacional impõe requisitos difíceis:

\begin{itemize}[leftmargin=*]
  \item eventos precisam de ordem e identidade estável;
  \item estado deve sobreviver a reinícios sem aceitar corrupção ou rollback;
  \item saída de modelo não pode virar verdade histórica nem conceder acesso;
  \item enforcement precisa ser determinístico, limitado e disponível durante
        falhas do plano de controle;
  \item toda ação negativa deve apontar para evidência durável e revisão exata
        de política;
  \item recuperação não pode ampliar uma negação para processos alheios.
\end{itemize}

O Unix-AGB atende a essas restrições separando observação, evidência canônica,
inferência de estado, compilação de política e enforcement no kernel. Essa
separação é mais importante que qualquer motor de estado específico.

Neste artigo, \emph{estado causal} significa estado dependente de trajetória
cuja interpretação de política depende de eventos predecessores ordenados no
mesmo namespace de segurança; não reivindicamos identificação causal no sentido
estatístico ou intervencionista. Portanto, o Unix-AGB não é primordialmente um
novo detector de trajetórias. É uma arquitetura para compilar com segurança
estado persistente e isolado por namespace em restrições de curta duração
voltadas ao kernel. O motor de estado determina quando uma trajetória justifica
restrição; o Unix-AGB determina como essa conclusão pode afetar o enforcement
sem dar ao motor autoridade para ampliar privilégios.

\subsection*{Contribuições}

Este trabalho apresenta quatro contribuições:

\begin{enumerate}[leftmargin=*]
  \item uma arquitetura de política para sistema operacional que separa estado
        causal persistente, evidência canônica e enforcement determinístico;
  \item compilação deny-only regida pelo Princípio da Autoridade Monotônica, no
        qual estado aprendido pode restringir, mas não ampliar privilégios;
  \item um pipeline de enforcement ciente de reinícios e revisões, com registros
        autenticados e expirantes de política e auditoria negativa durável; e
  \item um protótipo Linux experimental que demonstra enforcement seccomp com
        escopo de namespace, recuperação, rollback e semântica de falhas.
\end{enumerate}

\section{Trabalhos Relacionados e Posicionamento}

Linux Security Modules fornece uma infraestrutura comum no kernel para sistemas
de controle de acesso obrigatório como SELinux e AppArmor
~\cite{lsm,selinux,apparmor}. Seccomp user notification permite mediar em
userspace syscalls selecionadas~\cite{seccomp}, enquanto sistemas baseados em
eBPF aproximam observação programável e detecção de runtime dos eventos do
kernel~\cite{falco}. O Unix-AGB não substitui esses mecanismos: usa superfícies
existentes de enforcement Linux como última etapa de um pipeline de estado e
política versionado separadamente.

Trabalhos anteriores reconstruíram históricos de ataque a partir de proveniência
do sistema operacional~\cite{backtracker,camflow}, rastrearam fluxo de informação
e taint~\cite{denning}, correlacionaram streams com processamento de eventos
complexos e com estado~\cite{cep} e aprenderam comportamento anômalo de syscalls
~\cite{forrest}. EDRs comerciais baseados em comportamento seguem um padrão
operacional relacionado, embora sua semântica interna raramente constitua um
contrato público de sistemas. Esses trabalhos mostram que históricos podem
melhorar detecção e investigação. A questão distinta tratada aqui é como uma
conclusão persistente sobre trajetória se torna uma restrição de kernel compacta,
autenticada, expirante, vinculada a revisão, incapaz de conceder privilégio e
auditável através de reinício, rollback e falha de componentes. Qualidade de
detecção fica deliberadamente fora da contribuição central: FSMs, regras de
sequência, janelas, proxies e ASM-CM entram pela mesma fronteira de motor.

O ActPlane é especialmente próximo em motivação: conecta contexto de política
no nível do agente ao enforcement no sistema operacional, usa uma DSL de fluxo
de informação para políticas entre eventos e implementa mediação no kernel com
eBPF~\cite{actplane}. O Unix-AGB trata uma fronteira complementar de sistemas.
Em vez de permitir que um agente declare a política do kernel, considera o
estado de trajetória uma entrada derivada e potencialmente aprendida para um
compilador determinístico. Os registros resultantes são deny-only, autenticados,
vinculados a revisão, expirantes e ligados a evidência canônica durável. Sua
ênfase recai, portanto, em estado persistente através de reinícios e revisões de
política, autoridade monotônica, decisões negativas auditáveis e semânticas
explícitas de recuperação e falha.

\section{Decomposição do sistema}

A arquitetura é um pipeline de fronteiras de confiança estreitas:

\begin{center}
\begin{tabular}{>{\raggedright\arraybackslash}p{0.20\textwidth} >{\raggedright\arraybackslash}p{0.70\textwidth}}
\toprule
Camada & Responsabilidade \\
\midrule
Linux e coletores & Observar operações de processo, arquivo, identidade e rede; fornecer fatos estáveis do kernel e hooks de enforcement. \\
Event gateway & Validar schema, identidade, ordem, proveniência e tamanho antes de aceitar um evento. \\
Store canônico & Preservar fatos exatos em append-only; rejeitar duplicatas e replay de sequência. \\
Motor de estado & Manter estado limitado por namespace e emitir resumo com revisão, sinais, evidências e confiança opcional. \\
Motor de política & Aplicar invariantes determinísticas e traduzir estado validado em \texttt{ALLOW}, \texttt{DENY} ou \texttt{ABSTAIN}. \\
Cache compilado & Reter somente registros \texttt{DENY} autenticados, versionados e expirantes na borda de enforcement. \\
Adapter de enforcement & Vincular decisão compilada a um sujeito exato visível ao kernel e produzir resultado como \texttt{EACCES}. \\
Leitor opcional & Produzir explicações legíveis fora do caminho síncrono de enforcement. \\
\bottomrule
\end{tabular}
\end{center}

O store canônico e o motor de estado respondem a perguntas diferentes. O store
responde ``o que foi observado exatamente?''. O motor responde ``qual estado de
segurança limitado decorre do histórico aceito?''. Checkpoint, representação
neural ou acumulador de regras são artefatos derivados: podem ser descartados e
reconstruídos. A evidência canônica nunca deve ser reescrita silenciosamente para
concordar com eles.

\section{Eventos canônicos e identidade estável}

\subsection{Contrato de evento}

Cada evento aceito contém versão de schema, identificador global único,
sequência por namespace, tempos de parede e monotônico, host, namespace,
sujeito, operação, recurso, resultado solicitado ou observado, revisão de
política, rótulos e proveniência. O vocabulário inicial contém
\texttt{process.exec}, \texttt{process.exit}, \texttt{file.open},
\texttt{network.connect} e \texttt{identity.change}.

O contrato privilegia metadata. Ele não exige conteúdo de arquivo, payload de
pacote, segredo ou prompt. Resolver conteúdo seria uma operação autorizada
separada. Isso reduz exposição de privacidade e material controlado por atacante
no plano de política.

\subsection{Reuso de PID é um erro de identidade}

PID é reciclável. Estado indexado apenas por PID pode transferir o histórico de
um processo encerrado para outro. O Unix-AGB usa:

\begin{center}
\texttt{process:<boot-id>:<pid>:<start-time-ns>}
\end{center}

O boot ID separa reinicializações; o tempo de início separa reutilizações dentro
do mesmo boot. O gateway recalcula o namespace esperado e rejeita divergências.
No handoff do Gate 4, launcher e guardian derivam a identidade
independentemente de \texttt{/proc}; a mensagem é recusada se o namespace
declarado não corresponder ao processo vivo.

\subsection{Ordenação e lacunas}

A sequência cresce estritamente dentro de cada namespace. IDs duplicados e
sequências não crescentes são rejeitados antes de alterar o índice em memória.
Uma lacuna não é tratada como ausência benigna: torna o estado incompleto até
uma política explícita de ressincronização. Tempo de parede é descritivo;
sequência e tempo monotônico determinam a ordem causal.

Essa regra impede uma falha comum: interpretar perda de telemetria como prova de
que nada ocorreu.

\section{Estado causal persistente}

\subsection{Transição isolada por namespace}

O motor consome um evento validado por vez e emite
\texttt{SecurityStateSummary}. Sua revisão deve corresponder à sequência aceita.
O resumo contém faixa de risco, sinais, IDs de evidência, nome do motor,
fingerprint do checkpoint quando aplicável e horário. Consulta e atualização
exigem namespace exato; estado não é compartilhado implicitamente entre
processos, usuários, containers, serviços, agentes ou tenants.

Um exemplo determinístico mínimo memoriza uma execução, depois uma conexão
externa e finalmente um acesso a credencial. O último evento se torna elevado
somente quando as relações anteriores existem no mesmo namespace. A propriedade
central é o contrato de transição:

\[
S_{n+1} = F(S_n, E_{n+1}, C, R)
\]

onde $S_n$ é o estado do namespace, $E_{n+1}$ é o próximo evento canônico, $C$
é o fingerprint de configuração e $R$ é a revisão do motor ou checkpoint. Todo
consumidor deve identificar os quatro elementos.

\subsection{Motores substituíveis}

O Unix-AGB expõe uma fronteira de motor de estado em vez de acoplar política a
um modelo. Máquinas de estado, regras de sequência, janelas, proxy persistente e
adapters ASM-CM podem ser comparados pelos mesmos contratos. Isso viabiliza
baselines e impede que inferência aprendida redefina o protocolo de política.

O histórico experimental importa. Ensaios iniciais de generalização causal
foram negativos; identificadores canônicos locais à trajetória corrigiram uma
falha de representação e um ensemble congelado de três seeds passou num teste
sintético separado. Isso demonstra mecanismo e família controlada, não detecção
de ataques naturais desconhecidos. Por isso, divergência entre modelos resulta
por padrão em \texttt{ABSTAIN}, e inferência neural fica fora do hot path.

\subsection{Snapshots atômicos e restore}

Estado persistente só é útil se o restore for rigoroso. O protocolo vincula:

\begin{itemize}[leftmargin=*]
  \item versão do formato;
  \item fingerprint do motor;
  \item fingerprint da configuração;
  \item estado completo e ordenado dos namespaces;
  \item checksum ou autenticação.
\end{itemize}

A gravação usa arquivo temporário no mesmo diretório, flush, \texttt{fsync},
rename atômico e \texttt{fsync} do diretório. Se a persistência falha, a
atualização em memória é revertida. Restore rejeita snapshot corrompido,
incompatível ou de outra configuração em vez de inventar estado permissivo.

Checksums detectam corrupção acidental. Snapshots de política adicionam HMAC,
pois quem consegue editar o estado também consegue recalcular um hash comum.

\section{Do estado à política}

\subsection{Álgebra pequena de enforcement}

O core usa somente três efeitos:

\begin{description}[leftmargin=2.4cm,style=nextline]
  \item[\texttt{ALLOW}] Nenhum invariante restritivo foi encontrado. No Gate 3
        atual é resultado de shadow/auditoria, não permissão compilada.
  \item[\texttt{DENY}] Estado e política validados justificam restrição.
  \item[\texttt{ABSTAIN}] Evidência, confiança, revisão, persistência ou
        integridade contratual não bastam para decisão confiável.
\end{description}

Avisos são metadata de auditoria, não uma quarta ação do kernel. Assim, a
álgebra permanece pequena e avaliações de domínio passam por adapter antes da
camada de enforcement.

\subsection{Ordem da decisão e autoridade monotônica}

O Gate 3 valida configuração, contratos, igualdade de namespace, revisão de
política e revisão de estado. Estado \texttt{restricted} ou
\texttt{quarantined} produz \texttt{DENY}. Estado aprendido elevado exige
evidências, confiança mínima e fingerprint do checkpoint. Estado normal ou de
monitoramento produz \texttt{ALLOW} em shadow. Entrada desconhecida,
desatualizada, incompleta ou inconsistente produz \texttt{ABSTAIN} fail-closed.

Somente \texttt{DENY} é compilado. Nem o modelo nem cache ausente fabricam
permissão.

\paragraph{Princípio da Autoridade Monotônica.}
Decisões derivadas de estado podem preservar ou reduzir a autoridade concedida
pela política-base, mas nunca ampliá-la. A política Linux de base permanece,
portanto, autoritativa quando a inferência se abstém, o estado expira ou o plano
de controle com estado fica indisponível. A compilação deny-only é a regra
concreta de enforcement que realiza esse princípio no Unix-AGB.

\subsection{Durabilidade antes da restrição}

Uma decisão negativa é anexada à auditoria e sincronizada antes de entrar no
cache. Se a persistência falha, o efeito vira \texttt{ABSTAIN} com reason code
explícito. Registros positivos ou de abstention podem usar sincronização em
grupo, mas \texttt{DENY} força durabilidade imediata.

Isso estabelece um invariante recuperável:

\begin{center}
\textit{se o enforcement observa uma negação compilada, a evidência de auditoria
dessa negação já está durável.}
\end{center}

\section{Compilação autenticada de decisões}

Uma decisão compilada vincula namespace completo, operação, SHA-256 do recurso
canônico, revisões de política e estado, digest das evidências, ID e expiração.
O cache aceita somente \texttt{DENY}. Seu snapshot autentica registros ordenados
e revisão ativa com HMAC-SHA-256.

Lookup retorna \texttt{ABSTAIN} em miss, expiração, revisão incompatível ou
estado obsoleto. Rotação e rollback limpam estado incompatível. Um snapshot
antigo não pode ser renomeado como política nova. A projeção de runtime do Gate
4 ainda retém o último snapshot autenticado quando um reload é inválido; política
inicial inválida impede ativação, enquanto rotação vazia autenticada ou expiração
remove a restrição.

O protótipo atual de supervisão de serviço converte qualquer \texttt{DENY} ativo
no nível da trajetória em contenção de egress externo para o namespace exato.
A negação pode se originar numa leitura de credencial, escrita de persistência,
execução administrativa ou evento de rede; a operação originária permanece na
auditoria. Esse mapeamento é amplo e não deve ser chamado de
destination-specific: o runtime verifica namespace, revisão e expiração, mas
não compara o destino interceptado ao digest de recurso compilado.

\section{A borda de enforcement do Gate 4}

\subsection{Intercepção antes do exec}

O launcher opt-in cria o processo protegido e instala um filtro seccomp com user
notification antes de liberá-lo para \texttt{exec}. Isso fecha a corrida em que
a aplicação conectaria antes da proteção. O descritor do listener e um adapter
limitado para leitura de memória são transferidos ao guardian persistente por
socket Unix.

O handoff é autenticado por credenciais do peer, HMAC, revisão, nonce e expiração
monotônica curta. Descritores passam via \texttt{SCM\_RIGHTS}. O guardian valida
grupo de processos e recalcula o namespace. Replay, peer errado, revisão
obsoleta, expiração, namespace divergente ou transferência incompleta resultam
em rejeição.

\subsection{Decisão em \texttt{connect}}

Para cada notificação, o guardian:

\begin{enumerate}[leftmargin=*]
  \item valida que a notificação continua viva;
  \item resolve a thread para o processo do thread group;
  \item confere o escopo exato;
  \item lê no máximo 128 bytes do \texttt{sockaddr};
  \item permite loopback e sockets Unix conforme a política congelada;
  \item consulta o cache Gate 3 autenticado para egress externo;
  \item grava e sincroniza o registro de enforcement;
  \item retorna \texttt{EACCES} para negação correspondente ou continua a
        syscall nos demais casos.
\end{enumerate}

Família desconhecida, tamanho inválido, falha de leitura ou falha do adapter
tomam caminho fail-closed limitado ao alvo protegido. Processos fora do escopo
não herdam a negação.

\subsection{Propriedade do listener e recuperação}

Experimentos iniciais revelaram a diferença entre queda do worker de política e
perda do listener. Se o dono do listener seccomp morre, a syscall não recebe
necessariamente \texttt{EACCES}; pode ficar bloqueada ou seguir semântica de
teardown do kernel. Chamar isso de ``fail closed'' esconderia falha de
disponibilidade.

A solução mantém o listener no guardian e torna o worker substituível. IDs de
notificação são validados antes da resposta; falhas consomem orçamento limitado
de restart e usam fallback com escopo. Deadlines antes do lease e teardown do
processo protegido limitam casos em que ninguém assumiu a notificação com
segurança. A recuperação nunca amplia negação para outro processo.

\section{Empacotamento e operação}

O runtime de laboratório é um pacote Debian opt-in com unidade systemd, conta
dedicada, configuração em \texttt{/etc/unix-agb}, sockets em
\texttt{/run/unix-agb}, estado e auditoria em \texttt{/var/lib/unix-agb} e
chaves controladas por root. A configuração é vinculada a revisão e desabilitada
por padrão. A unidade usa hardening como \texttt{NoNewPrivileges},
\texttt{ProtectSystem=strict}, temporário privado, bounding de capabilities e
caminhos de escrita explícitos.

Testes cobrem instalação, default inativo, ativação, restart, daemon reload,
reboot, upgrade, disable, purge, remoção de conta, sockets e resíduos. Testes
negativos encontraram uma corrida: systemd indicava serviço simples ativo antes
do socket de controle existir. O launcher agora espera por intervalo limitado e
falha antes de liberar o filho se readiness não chegar.

Isso mostra por que packaging faz parte da segurança. Lógica correta no processo
não garante ordem de boot, upgrade, propriedade de segredo ou recuperação.

\section{Evidência experimental}

O repositório preserva protocolos, tentativas negativas, fingerprints e
relatórios, em vez de publicar somente sucessos. A progressão inclui:

\begin{itemize}[leftmargin=*]
  \item prova causal com operação terminal idêntica e decisões shadow diferentes;
  \item Gate 3 dry-run com revisões, auditoria durável, cache deny-only
        autenticado, expiração, corrupção e isolamento;
  \item piloto seccomp user notification retornando \texttt{EACCES};
  \item concorrência, overload, timeout e o resultado negativo de perda do
        listener;
  \item guardian com listener retido e recuperação do worker;
  \item enforcement empacotado e opt-in em VMs Ubuntu 24.04 e 26.04;
  \item transição de serviço persistente entre política vazia autenticada,
        negação no namespace exato, retenção após reload corrompido e retorno a
        política vazia autenticada.
\end{itemize}

As amostras variam de dezenas de microssegundos no guardian a poucos
milissegundos em brokers concorrentes. Elas caracterizam fronteiras controladas,
não capacidade de produção. Na supervisão, a aplicação protegida mudou de uma
tentativa externa permitida para \texttt{EACCES}, reteve negação após cache
malformado e recuperou após rotação vazia autenticada; o controle permaneceu
inalterado.

\begin{center}
\small
\begin{tabular}{>{\raggedright\arraybackslash}p{0.43\textwidth} >{\raggedright\arraybackslash}p{0.47\textwidth}}
\toprule
Experimento & Resultado observado \\
\midrule
Guardian com listener retido, 256 chamadas-alvo~\cite{guardian} & 256/256 retornaram \texttt{EACCES}; latência de decisão de 44~$\mu$s p50, 49~$\mu$s p95 e 68~$\mu$s p99 \\
Broker concorrente, 256 chamadas-alvo~\cite{concurrency} & 256/256 retornaram \texttt{EACCES}; 1,148~ms p50, 1,939~ms p95 e 2,366~ms p99 \\
Falha de worker em voo~\cite{inflight} & A notificação pendente exata recebeu \texttt{EACCES} em 123~$\mu$s; worker substituto pronto em 9,839~ms \\
Enforcement exact-launch em VM~\cite{exactlaunch} & Conexão externa protegida recebeu \texttt{EACCES}; loopback e controle não protegido permaneceram inalterados \\
Transições do cache autenticado & Reload corrompido reteve a última negação válida; rotação vazia autenticada a removeu \\
Auditoria na supervisão de serviço & 191~$\mu$s p50 e 276~$\mu$s p95/p99, incluindo auditoria durável \\
\bottomrule
\end{tabular}
\end{center}

A tabela relata fronteiras experimentais distintas e não deve ser interpretada
como uma única distribuição de latência end-to-end. Em particular, o
microbenchmark do guardian, o broker concorrente, as VMs empacotadas e a
supervisão usaram workloads controlados diferentes.

Uma transição posterior substituiu o trigger artesanal pelo ensemble congelado
de três membros e pelo compilador real de política Gate 3 enquanto o serviço
protegido permanecia vivo. O \texttt{DENY} autenticado, originado em
\texttt{file.open}, fez a conexão externa seguinte receber \texttt{EACCES}.
Entretanto, a entrada foi uma trajetória controlada congelada e explicitamente
rebound ao namespace vivo, não telemetria BPF recém-capturada. A matriz
congelada mantém o Gate 4 como protótipo controlado.

\section{Semântica de falha}

``Fail closed'' precisa nomear escopo e consequência. O Unix-AGB distingue:

\begin{itemize}[leftmargin=*]
  \item \textbf{abstention de política}: não há decisão confiável derivada de estado;
  \item \textbf{negação com escopo}: o request protegido recebe \texttt{EACCES};
  \item \textbf{teardown com escopo}: o workload é encerrado quando resposta
        segura não pode ser garantida;
  \item \textbf{continuidade da política-base}: Linux continua ativo se modelo
        ou plano de controle falha;
  \item \textbf{falha de disponibilidade}: request trava ou perde resposta útil,
        sem ser indevidamente declarado como negação bem-sucedida.
\end{itemize}

Esse vocabulário obriga o projeto a tratar separadamente telemetria ausente,
lacunas, snapshot corrompido, decisão expirada, queda de worker, perda de
listener, falha de auditoria, rollback e corrida de restart.

\section{Propriedades e riscos residuais}

Os contratos implementados sustentam propriedades concretas:

\begin{itemize}[leftmargin=*]
  \item identidade estável resiste a reuso de PID;
  \item namespace exato limita vazamento de estado entre processos;
  \item replay e sequência preservam ordem;
  \item cache autenticado e revision-bound resiste a substituição não autorizada;
  \item compilação deny-only impede ampliação de privilégio por modelo;
  \item expiração e rotação vazia autenticada fornecem rollback limitado;
  \item auditoria durável precede enforcement negativo;
  \item verificação independente vincula listener ao alvo vivo.
\end{itemize}

Os riscos residuais são relevantes. Coletor privilegiado comprometido pode
forjar observações; kernel comprometido derrota observação e enforcement;
chaves HMAC precisam de provisionamento real; binding por destino está
incompleto na projeção atual; attachment sistêmico não cooperativo não existe;
eficácia em telemetria natural não foi provada; e fail-closed amplo pode virar
negação de serviço. São bloqueadores de promoção, não detalhes editoriais.

\section{Por que o estado causal fica ao lado do kernel}

O protótipo ainda não justifica kernel próprio. Linux já fornece as superfícies
experimentais: observação por BPF/audit, identidade de processo, credenciais de
peer, cgroups, systemd, seccomp user notification e LSMs. Estado causal,
checkpoints, revisão, auditoria e inferência opcional evoluem mais rápido e têm
superfície de falha maior que um hook mínimo.

Mantê-los em userspace permite restart, substituição, comparação com baselines e
abstention explícita. A borda recebe uma restrição compacta, autenticada e
expirante, não um modelo. Trabalho de kernel só deve ser considerado se medições
revelarem hook ausente, identidade indisponível, latência inaceitável ou
propriedade impossível nas interfaces atuais.

\section{Conclusão}

Estado causal persistente pode enriquecer enforcement somente quando projetado
como pipeline restrito de evidência e política, não como modelo opaco no hot
path. O padrão do Unix-AGB é:

\begin{center}
eventos exatos $\rightarrow$ estado persistente isolado $\rightarrow$ resumo
validado $\rightarrow$ política determinística $\rightarrow$ negação autenticada
e expirante $\rightarrow$ enforcement com escopo.
\end{center}

As fronteiras duráveis são a contribuição: fatos separados de estado derivado;
identidade resistente a PID reuse; incerteza como \texttt{ABSTAIN}; apenas
\texttt{DENY} compilado; auditoria antes da restrição; listener disponível
durante recuperação da política. Estado Causal Persistente, Compilação Deny-Only
e Autoridade Monotônica pretendem ser conceitos aplicáveis além deste runtime
específico. A evidência atual mostra que os mecanismos podem ser montados e
testados num protótipo Linux empacotado e opt-in. Não mostra
prontidão de produção nem detecção geral de ataques desconhecidos. A promoção
exige a matriz congelada: binding independente e vivo de decisão Gate 3,
isolamento concorrente, endurance, recursos, update/rollback autenticado,
injeção de falhas, aplicações reais e lifecycle multi-VM sob uma única revisão
e artefato.

\begin{thebibliography}{99}

\bibitem{lsm}
C. Wright et al.,
\emph{Linux Security Modules: General Security Support for the Linux Kernel},
11th USENIX Security Symposium, 2002.

\bibitem{selinux}
P. Loscocco e S. Smalley,
\emph{Integrating Flexible Support for Security Policies into the Linux
Operating System}, USENIX Annual Technical Conference, 2001.

\bibitem{apparmor}
Linux Kernel Documentation,
\emph{AppArmor}, \url{https://docs.kernel.org/admin-guide/LSM/apparmor.html}.

\bibitem{seccomp}
Linux Kernel Documentation,
\emph{Seccomp BPF User Notification},
\url{https://docs.kernel.org/userspace-api/seccomp_filter.html}.

\bibitem{falco}
The Falco Authors,
\emph{Falco: Cloud Native Runtime Security}, \url{https://falco.org/docs/}.

\bibitem{backtracker}
S. T. King e P. M. Chen,
\emph{Backtracking Intrusions}, 19th ACM Symposium on Operating Systems
Principles, 2003.

\bibitem{camflow}
T. Pasquier et al.,
\emph{Practical Whole-System Provenance Capture}, ACM Symposium on Cloud
Computing, 2017.

\bibitem{denning}
D. E. Denning,
\emph{A Lattice Model of Secure Information Flow}, Communications of the ACM,
19(5), 1976.

\bibitem{cep}
G. Cugola e A. Margara,
\emph{Processing Flows of Information: From Data Stream to Complex Event
Processing}, ACM Computing Surveys, 44(3), 2012.

\bibitem{forrest}
S. Forrest et al.,
\emph{A Sense of Self for Unix Processes}, IEEE Symposium on Security and
Privacy, 1996.

\bibitem{actplane}
Y. Zheng et al.,
\emph{ActPlane: Programmable OS-Level Policy Enforcement for Agent Harnesses},
arXiv:2606.25189v2, 2026,
\url{https://arxiv.org/abs/2606.25189v2}.

\bibitem{architecture}
Projeto Unix-AGB,
\emph{Especificação de Arquitetura Unix-AGB},
\texttt{docs/Unix\_AGB\_Architecture\_Specification\_PTBR.md}, 2026.

\bibitem{contracts}
Projeto Unix-AGB,
\emph{Contratos AGB},
\texttt{docs/contracts.md}, 2026.

\bibitem{gate3}
Projeto Unix-AGB,
\emph{Gate 3 Dry-Run Policy Engine},
\texttt{docs/gate3-policy-engine.md}, 2026.

\bibitem{promotion}
Projeto Unix-AGB,
\emph{Gate 4 Promotion Protocol},
\texttt{docs/report/106\_gate4\_promotion\_protocol\_20260817.md}, 2026.

\bibitem{supervision}
Projeto Unix-AGB,
\emph{Gate 4 Gate 3 Service Supervision},
\texttt{docs/report/108\_gate4\_gate3\_service\_supervision\_20260817.md}, 2026.

\bibitem{concurrency}
Projeto Unix-AGB,
\emph{Gate 4 Concurrent Seccomp Broker},
\texttt{docs/report/94\_gate4\_concurrent\_seccomp\_broker\_20260817.md}, 2026.

\bibitem{guardian}
Projeto Unix-AGB,
\emph{Gate 4 Retained-Listener Guardian},
\texttt{docs/report/96\_gate4\_retained\_listener\_guardian\_20260817.md}, 2026.

\bibitem{inflight}
Projeto Unix-AGB,
\emph{Gate 4 In-Flight Notification Recovery},
\texttt{docs/report/97\_gate4\_inflight\_notification\_recovery\_20260817.md}, 2026.

\bibitem{exactlaunch}
Projeto Unix-AGB,
\emph{Gate 4 Persistent Exact-Launch Enforcement},
\texttt{docs/report/105\_gate4\_persistent\_exact\_launch\_enforcement\_20260817.md}, 2026.

\end{thebibliography}

\end{document}
